\documentclass[12pt]{article}
\usepackage[utf8]{inputenc}
\usepackage[T1]{fontenc}
\usepackage[margin=1in]{geometry}
\usepackage{setspace}
\usepackage{xcolor}
\usepackage[colorlinks=true, urlcolor=blue!60!black, linkcolor=blue!60!black]{hyperref}
\setlength{\parindent}{0em}
\setlength{\parskip}{0.45em}
\renewcommand{\baselinestretch}{1.0}
\begin{document}
\begin{center}
{\Large \textit{303 Creative LLC v.\ Elenis}, 600 U.S.\ 570 (2023)\par}
\vspace{0.3em}
{\normalsize Case Brief --- One-Page Edition\par}
\end{center}

\textbf{Memory Jogger.} A state may not use its public accommodations law to compel a website designer to create custom wedding websites conveying messages she rejects (\href{https://supreme.justia.com/cases/federal/us/600/21-476/}{Justia}).

\textbf{Facts.} Lorie Smith of 303 Creative LLC wanted to design wedding websites only for opposite-sex weddings, citing her Christian beliefs. Colorado's Anti-Discrimination Act bars public accommodations from refusing service based on sexual orientation. The parties stipulated that her websites are ``original, customized'' and ``expressive'' and that she serves LGBTQ+ clients in other work.

\textbf{Procedural History.} In a pre-enforcement suit, the district court ruled for Colorado; the Tenth Circuit affirmed 2--1, finding CADA survived strict scrutiny.

\textbf{Votes.} Reversed, 6--3: Gorsuch for the Court; Sotomayor, with Kagan and Jackson, dissented.

\textbf{Holding.} A state may not compel a business that creates customized, expressive ``pure speech'' to convey messages it disagrees with, even when it sells to the public for profit.

\textbf{Reasoning.} Under \textit{Barnette}, \textit{Hurley}, and \textit{Dale}, government may not compel a speaker's message. Smith's websites are her own speech; profit and corporate form do not matter; her objection is to a message, not a status; and unique talent cannot justify coercion. Sotomayor answered that CADA regulates the conduct of refusing service, with only an ``incidental'' burden on speech, and that the Court had never before granted businesses a right to discriminate.

\textbf{Examples.} \emph{Same-Side:} a custom wedding photographer. \emph{Opposite-Side:} a reception venue renting space and serving food. \emph{Fence-Sitter:} a custom cake baker or florist.

\textbf{Critique.} Scholars fault the opinion for naming no tier of scrutiny or attribution test, and for resting on an abstract, stipulated record.

\textbf{Key Quotation.} ``When a state public accommodations law and the Constitution collide, there can be no question which must prevail.'' --- Justice Gorsuch, for the Court.

\end{document}
